% =====================================================================
% 独立冊子『逆引きチートシート』本文
% 元原稿: 05_逆引きチートシート.md（全範囲）＋ 02_GitHub編.md（SSH/PAT）
% =====================================================================
\chapter{逆引きチートシート}

本編『はじめてのバージョン管理』と別冊『チームで安全に開発する』で登場したコマンドを \textbf{「〜したい → このコマンド」} の形で引ける辞書です。あわせて、本編では扱わなかった\textbf{SSH鍵・パーソナルアクセストークン（PAT）}による認証の設定方法も、本書の最後にまとめます。使い方や背景を知りたいときは、各節末尾に示した参照先を見てください。

\begin{hintbox}
\textbf{迷ったらまず \cmd{git status}}：今どのファイルがどのエリアにあるかが分かり、次の一手が見えます。
\end{hintbox}

% ---------------------------------------------------------------------
\section{最初の設定（1度だけ）}

\begin{center}
\begin{tabularx}{\linewidth}{@{}X X@{}}
\toprule
したいこと & コマンド \\
\midrule
名前を登録する & \cmd{git config -{}-global user.name "your-name"} \\
メールアドレスを登録する & \cmd{git config -{}-global user.email "you@example.com"} \\
設定内容を確認する & \cmd{git config -{}-global -{}-list} \\
新規リポジトリの初期ブランチ名を \cmd{main} にする & \cmd{git config -{}-global init.defaultBranch main} \\
Gitのバージョンを確認する & \cmd{git -{}-version} \\
\bottomrule
\end{tabularx}
\end{center}

詳しくは本編「最初の設定」を参照。

% ---------------------------------------------------------------------
\section{変更を記録する（黄金フロー）}

\textbf{編集 → \cmd{add} → \cmd{commit}} の3ステップが基本です。

\begin{center}
\begin{tabularx}{\linewidth}{@{}X X@{}}
\toprule
したいこと & コマンド \\
\midrule
新しいフォルダをリポジトリにする & \cmd{git init} \\
変更した1ファイルをステージングに追加する & \cmd{git add ファイル名} \\
変更した全ファイルをステージングに追加する & \cmd{git add .} \\
コミット（セーブポイント）を作る & \cmd{git commit -m "変更内容のメモ"} \\
ステージング＋コミットをまとめて行う（追跡済みのみ） & \cmd{git commit -am "メモ"} \\
\bottomrule
\end{tabularx}
\end{center}

詳しくは本編「はじめてのコミット」を参照。

% ---------------------------------------------------------------------
\section{状態・履歴を確認する}

\begin{center}
\begin{tabularx}{\linewidth}{@{}X X@{}}
\toprule
したいこと & コマンド \\
\midrule
どのファイルが変更され、どのエリアにあるか見る & \cmd{git status} \\
コミット履歴を見る & \cmd{git log} \\
履歴を1行ずつ簡潔に見る & \cmd{git log -{}-oneline} \\
まだステージングしていない変更の中身を見る & \cmd{git diff} \\
ステージング済みの変更の中身を見る & \cmd{git diff -{}-staged} \\
リモートの登録内容を確認する & \cmd{git remote -v} \\
いまのブランチの「送り先」を確認する & \cmd{git branch -vv} \\
\bottomrule
\end{tabularx}
\end{center}

詳しくは本編「状態を確認するコマンド」を参照（\cmd{git log} の表示は \cmd{q} キーで終了）。

% ---------------------------------------------------------------------
\section{認証・リモートとつなぐ・同期する}

\cmd{push} する前に、GitHub CLI（\cmd{gh}）で1度だけログインします。

\begin{center}
\begin{tabularx}{\linewidth}{@{}X X@{}}
\toprule
したいこと & コマンド \\
\midrule
\cmd{gh} がインストール済みか確認する & \cmd{gh -{}-version} \\
GitHubにログインする & \cmd{gh auth login} \\
ログイン状態を確認する & \cmd{gh auth status} \\
GitHubのリポジトリを丸ごと複製する & \cmd{git clone URL} \\
手元のリポジトリにリモートを登録する & \cmd{git remote add origin URL} \\
リモートのURLを付け替える & \cmd{git remote set-url origin 新しいURL} \\
ブランチ名を \cmd{main} に統一する & \cmd{git branch -M main} \\
初回のアップロード（送り先を覚えさせる） & \cmd{git push -u origin main} \\
2回目以降のアップロード & \cmd{git push} \\
リモートの最新を取得して手元に反映する & \cmd{git pull} \\
リモートの最新を取得だけする（反映しない） & \cmd{git fetch} \\
\bottomrule
\end{tabularx}
\end{center}

詳しくは本編「GitHubにログインする」「ローカルとリモートをつなぐ」を参照。SSH鍵・PATを使う場合は本書末尾「認証の別方式」へ。

\begin{hintbox}
\textbf{始めに \cmd{pull}、終わりに \cmd{push}} を習慣にすると、トラブルが起きにくくなります。
\end{hintbox}

% ---------------------------------------------------------------------
\section{ブランチを操作する}

\begin{center}
\begin{tabularx}{\linewidth}{@{}X X@{}}
\toprule
したいこと & コマンド \\
\midrule
ブランチ一覧を見る（\cmd{*} が今いるブランチ） & \cmd{git branch} \\
ブランチを作って同時に切り替える & \cmd{git switch -c feature/xxx} \\
（古い書き方）作って切り替える & \cmd{git checkout -b feature/xxx} \\
既存のブランチへ切り替える & \cmd{git switch ブランチ名} \\
作業ブランチをリモートへ上げる（初回） & \cmd{git push -u origin feature/xxx} \\
ローカルのブランチを削除する（マージ済みのみ安全に） & \cmd{git branch -d feature/xxx} \\
リモートのブランチを削除する & \cmd{git push origin -{}-delete feature/xxx} \\
\bottomrule
\end{tabularx}
\end{center}

詳しくは別冊『チームで安全に開発する』の「ブランチの基本操作」を参照。

\begin{warnbox}
作業を始める前に、\textbf{\cmd{main} を最新にしてから枝分かれ}するのが安全です（\cmd{git switch main} → \cmd{git pull} → \cmd{git switch -c 新ブランチ}）。
\end{warnbox}

% ---------------------------------------------------------------------
\section{PRとマージ後の後片付け・コンフリクト解決}

PRの作成・レビュー・マージは\textbf{GitHubの画面（ブラウザ）で操作}します。後片付けとコンフリクト解決はコマンドで行います。

\begin{center}
\begin{tabularx}{\linewidth}{@{}X X@{}}
\toprule
したいこと & コマンド／操作 \\
\midrule
PRを作る & GitHubの「Compare \& pull request」→ base/compare を確認 → 作成 \\
PRに変更を追加する & 同じブランチで \cmd{add} → \cmd{commit} → \cmd{push} \\
マージ後にローカルの \cmd{main} を最新化する & \cmd{git switch main} → \cmd{git pull} \\
マージ済みのローカルブランチを消す & \cmd{git branch -d feature/xxx} \\
マージ済みのリモートブランチを消す & \cmd{git push origin -{}-delete feature/xxx} \\
どのファイルが競合したか確認する & \cmd{git status} \\
解決したファイルを「解決済み」として登録する & \cmd{git add ファイル名} \\
解決を確定する（メッセージは自動用意） & \cmd{git commit} \\
マージをやめて開始前の状態に戻す & \cmd{git merge -{}-abort} \\
\bottomrule
\end{tabularx}
\end{center}

詳しくは別冊『チームで安全に開発する』の該当章（「プルリクエスト（PR）を作る」・「コンフリクト（競合）を解決する」）を参照。

% ---------------------------------------------------------------------
\section{やり直し・取り消し}

\begin{warnbox}
取り消し系は状態を戻す操作です。\textbf{不安なときはまず \cmd{git status} で現在地を確認}してから使ってください。
\end{warnbox}

\begin{center}
\begin{tabularx}{\linewidth}{@{}X X@{}}
\toprule
したいこと & コマンド \\
\midrule
作業ツリーの変更を捨てて最後のコミットに戻す（1ファイル） & \cmd{git restore ファイル名} \\
ステージングから下ろす（変更内容は残す） & \cmd{git restore -{}-staged ファイル名} \\
直前のコミットメッセージを修正する & \cmd{git commit -{}-amend -m "新しいメモ"} \\
マージ作業を中止して開始前に戻す & \cmd{git merge -{}-abort} \\
公開済みのコミットを打ち消す（履歴は残す） & \cmd{git revert コミットID} \\
\bottomrule
\end{tabularx}
\end{center}

各操作の使い分けや事故からの復旧は『トラブルシューティング』を参照。

% ---------------------------------------------------------------------
\section{困ったときのフローチャート}

「今こうしたい」から、打つべきコマンドを引けます（図\ref{fig:help-flow}）。

\begin{numfig}{「〜したい」から引く早見表}{fig:help-flow}
\begin{tikzpicture}[node distance=0.55cm and 1.4cm]
  \node[decisionnode, fill=gitAmber, text width=3cm] (Q1) {今の状態が分からない};
  \node[flownode, fill=gitOrange, text width=6.5cm, right=of Q1] (A1) {git status};
  \node[decisionnode, fill=gitAmber, text width=3cm, below=of Q1] (Q2) {手元の変更を上げたい};
  \node[flownode, fill=gitGreen, text width=6.5cm, right=of Q2] (A2) {git add . → git commit → git push};
  \node[decisionnode, fill=gitAmber, text width=3cm, below=of Q2] (Q3) {クラウドの最新を取り込みたい};
  \node[flownode, fill=gitBlue, text width=6.5cm, right=of Q3] (A3) {git pull};
  \node[decisionnode, fill=gitAmber, text width=3cm, below=of Q3] (Q4) {新しい作業を始めたい};
  \node[flownode, fill=gitPurple, text width=6.5cm, right=of Q4] (A4) {git switch main → git pull → git switch -c feature/xxx};
  \node[decisionnode, fill=gitAmber, text width=3cm, below=of Q4] (Q5) {pull / merge で競合が出た};
  \node[flownode, fill=gitRed, text width=6.5cm, right=of Q5] (A5) {マーカーを消す → git add → git commit};
  \draw[flowarrow] (Q1) -- (A1);
  \draw[flowarrow] (Q2) -- (A2);
  \draw[flowarrow] (Q3) -- (A3);
  \draw[flowarrow] (Q4) -- (A4);
  \draw[flowarrow] (Q5) -- (A5);
\end{tikzpicture}
\end{numfig}

各コマンドの「なぜそうするのか」は、本編・別冊の本文で図とともに説明しています。手が止まったら、まずこのチートシートで当たりをつけ、詳しく知りたくなったら該当箇所へ戻ってください。

% ---------------------------------------------------------------------
\section{認証の別方式（SSH鍵・PAT）}

本編では、最も手軽な \textbf{GitHub CLI（\cmd{gh}）によるHTTPS認証}だけを扱いました。ここでは、頻繁に使う人向けの \textbf{SSH鍵}と、トークンを手動管理する \textbf{PAT（パーソナルアクセストークン）}の設定方法を補足します。\textbf{どれか1つを設定すればOK}です。

\begin{center}
\begin{tabularx}{\linewidth}{@{}l l X@{}}
\toprule
方式 & URLの形 & 特徴 \\
\midrule
\textbf{HTTPS（gh）} & \cmd{https://github.com/\ldots.git} & \textbf{本編で採用}。\cmd{gh auth login} だけで使えて最も手軽 \\
\textbf{SSH鍵} & \cmd{git@github.com:\ldots.git} & 鍵を1度登録すれば以降は認証不要。複数リポジトリで使い回せる \\
\textbf{PAT} & \cmd{https://github.com/\ldots.git} & トークンを手動で発行・管理したい場合 \\
\bottomrule
\end{tabularx}
\end{center}

\subsection{SSH鍵}

SSHは、\textbf{鍵ペア（秘密鍵と公開鍵）}を使う方式です。公開鍵をGitHubに登録しておくと、以降は認証情報の入力なしで \cmd{push} できます。

\noindent\textbf{手順1：鍵ペアを作る}（メールアドレスはGitHubのもの）

\begin{codebox}[bash]
\begin{code}
ssh-keygen -t ed25519 -C "you@example.com"
\end{code}
\end{codebox}

\begin{itemize}[leftmargin=1.4em]
  \item 保存先を聞かれたら、そのまま \textbf{Enter}（デフォルトの \verb|~/.ssh/id_ed25519| に保存されます）
  \item パスフレーズは、設定してもしなくても構いません（設定するとより安全）
\end{itemize}

これで秘密鍵 \verb|id_ed25519| と公開鍵 \verb|id_ed25519.pub| が \verb|~/.ssh/| に作られます。

\begin{warnbox}
\textbf{秘密鍵（\cmd{.pub} が付かない方）は絶対に他人に渡さないでください。} GitHubに登録するのは公開鍵（\cmd{.pub} が付く方）だけです。
\end{warnbox}

\noindent\textbf{手順2：公開鍵の中身をコピーする}

\begin{codebox}[bash]
\begin{code}
# macOS
pbcopy < ~/.ssh/id_ed25519.pub

# Windows（Git Bash）
cat ~/.ssh/id_ed25519.pub | clip

# Linux（表示された内容を手動でコピー）
cat ~/.ssh/id_ed25519.pub
\end{code}
\end{codebox}

\noindent\textbf{手順3：GitHubに公開鍵を登録する}

\begin{enumerate}[leftmargin=1.6em]
  \item GitHubの \textbf{Settings} → \textbf{SSH and GPG keys} を開きます。
  \item \textbf{New SSH key} をクリックします。
  \item \textbf{Title} に分かりやすい名前（例：\cmd{my-laptop}）、\textbf{Key} に手順2でコピーした内容を貼り付けます。
  \item \textbf{Add SSH key} で登録完了です。
\end{enumerate}

\screenshot[0.62]{02_github/ssh_key_add.png}{SSH and GPG keys の New SSH key 登録画面}

\noindent\textbf{手順4：接続を確認する}

\begin{codebox}[bash]
\begin{code}
ssh -T git@github.com
\end{code}
\end{codebox}

次のように表示されれば成功です（\cmd{yes} を求められたら \cmd{yes} と入力）。

\begin{codebox}[text]
\begin{code}
Hi your-github-name! You've successfully authenticated, but GitHub does not provide shell access.
\end{code}
\end{codebox}

\subsection{パーソナルアクセストークン（PAT）}

\cmd{gh} を使わずHTTPSで \cmd{push} すると、ユーザー名とパスワードを聞かれます。ここで\textbf{GitHubのログインパスワードは使えません}。代わりに \textbf{PAT} という専用の文字列を発行して使います。

\noindent\textbf{トークンの発行手順}

\begin{enumerate}[leftmargin=1.6em]
  \item GitHubにログインし、右上のアイコン → \textbf{Settings} を開きます。
  \item 左メニュー最下部の \textbf{Developer settings} をクリックします。
  \item \textbf{Personal access tokens} → \textbf{Tokens (classic)} → \textbf{Generate new token (classic)} を選びます。
  \item \textbf{Note}（用途のメモ）、\textbf{Expiration}（有効期限）を設定し、\textbf{Select scopes} で \cmd{repo} にチェックして発行します。
  \item \textbf{Generate token} を押すと、\verb|ghp_xxxx…| という文字列が表示されます。
\end{enumerate}

\screenshot[0.62]{02_github/pat_scopes.png}{トークン発行画面（repoスコープにチェック）}

\begin{warnbox}
\textbf{トークンは発行直後の1回しか表示されません。} その場でコピーし、パスワードマネージャ等に保管してください。画面を閉じると二度と確認できず、再発行が必要になります。
\end{warnbox}

\noindent\textbf{トークンの使い方}：\cmd{push} 時にユーザー名とパスワードを聞かれたら、\textbf{パスワードの欄に、先ほどのトークンを貼り付けます}。一度入力すれば記憶され、次回以降は不要になります。

\begin{codebox}[text]
\begin{code}
Username: your-github-name
Password: ghp_xxxxxxxxxxxxxxxxxxxx   ← ログインパスワードではなくトークンを入力
\end{code}
\end{codebox}
